\chapter*{ЦЕЛЬ И ЗАДАЧИ РАБОТЫ}
\addcontentsline{toc}{chapter}{ЦЕЛЬ И ЗАДАЧИ РАБОТЫ}

Цель работы --- изучить методы экспериментального исследования производительности центрального процессора и получить практические навыки анализа влияния подсистемы памяти и особенностей микроархитектуры CPU на выполнение программ.

Для достижения цели были поставлены и решены следующие задачи:
\begin{enumerate}
	\item определить характеристики используемой вычислительной системы и подготовить средства аппаратного профилирования;
	\item собрать и запустить тестовое приложение, выполнить базовое профилирование и определить основные показатели производительности CPU: время выполнения, \texttt{cycles}, \texttt{instructions} и IPC;
	\item исследовать повторяемость измерений производительности и оценить разброс результатов;
	\item исследовать влияние размера рабочего набора и поведения кэш-памяти на производительность программы;
	\item исследовать влияние шага доступа к памяти \texttt{stride} и пространственной локальности обращений;
	\item исследовать влияние типа данных на размер рабочего набора и характеристики выполнения программы;
	\item исследовать работу TLB и влияние обращений к различным страницам памяти;
	\item исследовать влияние предсказуемых и непредсказуемых условных переходов на производительность;
	\item проанализировать полученные результаты и связать наблюдаемые эффекты с характеристиками используемой микроархитектуры.
\end{enumerate}

Работа выполнялась на Apple M3 Pro под macOS, где \texttt{perf} недоступен (раздел~\ref{sec:tools}).
